\sscp{Binding: spreading of high vowels across morpheme boundaries}{13-02-03}
\citet[250--254]{Mills1991} sees the feature of what he calls “binding”
as forming the basis for a “Southeast Maluku” subgroup,
and also being in support of CMP and possibly CEMP.
He describes it as,

\begin{quote}
		“Binding is a morphophonemic process whereby at certain junctures
		(1) a final (i,u) of the first element
		and the initial C of the second element metathesise;
		(2) the shifted V reduces to (y, w);
		(3) the original V is deleted.
		(Before a V-initial base,
		the final (i,u) simply shifts to (y, w).)
		Thus, the two elements are in a sense interlaced,
		bound together into a single entity.
		Very schematically: {\#}{\ldots} (C) (i,u) + (C)V{\ldots} ===>
		{\#}{\ldots} (C){\0}(C) (y, w)V{\ldots}” \citep[250]{Mills1991}.
\end{quote}

This, and similar phenomena, are also described by various scholars
as “pseudometathesis”, “morpheme external metathesis”, “a sandhi phenomenon”,
and “spreading of features across morpheme boundaries”.

\citeauthor{Mills1991} finds this process of ``binding'' operating in subject prefixes
on verbs, in compounding, and with suffixes as illustrated in examples
\qf{13-13}--\qf{13-17} using examples from \citet{Mills1991},
repackaged and updated here).

\noindent{\st{6pt}\tblr
{>{\hspace{-\tabcolsep}}l<{\hspace{-\tabcolsep}}
>{\hspace{-\tabcolsep}\enspace}
lll}
\tbx{13-13} & \mc{3}{l}{Fordata subject prefixes} \\
						& \it{dava}		&						& `seek'\\
						&	\it{mdwava}	&	/mu-dava/	&	`\tsc{2sg}-seek'	\\
						&	\it{mdyava}	&	/mi-dava/	&	`\tsc{2pl}-seek'	\\
\end{tabular}\mbox{}\\}

\noindent{\st{6pt}\tblr
{>{\hspace{-\tabcolsep}}l<{\hspace{-\tabcolsep}}
>{\hspace{-\tabcolsep}\enspace}
ll}
\tbx{13-14} & \mc{2}{l}{Yamdena compounding} \\
						& \it{bali}		& `side'\\
						&	\it{duwe}		&	`two'	\\
						&	\it{baldyu}	&	`both sides'	\\
\end{tabular}\mbox{}\\}

\noindent{\st{6pt}\tblr
{>{\hspace{-\tabcolsep}}l<{\hspace{-\tabcolsep}}
>{\hspace{-\tabcolsep}\enspace}
lll}
\tbx{13-15} & \mc{3}{l}{Leti and Moa compounding \citep[172]{Jonker1932}} \\
						&				& \it{pipi}			& `goat'\\
						&				&	\it{duma}			&	< Malay \it{domba} `sheep'	\\
						&	Leti	&	\it{pipi-duma}	&	`sheep'	\\
						&	Moa		&	\it{pipdiuma}		&	`sheep'	\\
\end{tabular}\mbox{}\\}

\noindent{\st{6pt}\tblr
{>{\hspace{-\tabcolsep}}l<{\hspace{-\tabcolsep}}
>{\hspace{-\tabcolsep}\enspace}
lll@{ }l@{ }l@{ }l@{ }l}
\tbx{13-16} & \mc{7}{l}{Yamdena suffixing} \\
						&	\it{manik}		& `bird'	&&&&& \\
						&	\it{mankyar}	&	`birds'	&	/manik/ + /-ar/ &{\ra}& **manki-ar	&{\ra}& \it{manky-ar}	\\
						&	\it{lutur}		&	`wall'	&&&&& \\
						&	\it{lutrwar}	&	`walls'	&	/lutur/ + /-ar/ &{\ra}& **lutru-ar	&{\ra}& \it{lutrw-ar}	\\
\end{tabular}\mbox{}\\}

\noindent{\st{6pt}\tblr
{>{\hspace{-\tabcolsep}}l<{\hspace{-\tabcolsep}}
>{\hspace{-\tabcolsep}\enspace}
lll}
\tbx{13-17} & \mc{3}{l}{Selaru suffixing} \\
						& \it{hahy}		&	`pig'			& \\
						&	\it{hahkye}	&	`the pig'	&	/hahy/ + /-ke/	\\
						&	\it{hahire}	&	`pigs'		&	/hahy/ + /-Vre/	\\
\end{tabular}\mbox{}\\}

\citet{CowardCoward2000} point out that different behaviour in
the morpho-phonology requires that Selaru distinguish consonants
(C), vowels (V), and glides (G) as distinct emic categories in the phonotactics of Selaru.
The data and analysis in example \qf{13-18} were provided by
David Coward (pers. comm. 21 Sept.
2020), illustrating a fuller paradigm for Selaru.
For body parts, the clearest context to see whether a Selaru
root ends in a V or G is if it is dead or disembodied,
“or found in a cave or such (not possessed or associated with
some person or animal) [{\ldots}] But if it is alive/possessed,
then it surfaces with an inserted Vowel beat for \tsc{gen} affixation
(the vowel beat causes glides to become syllabic,
but inserts an \it{a} when there are no glides)”.

\noindent{\wg\st{6pt}\begin{tabularx}
{\linewidth}{>{\hspace{-\tabcolsep}}l<{\hspace{-\tabcolsep}}
>{\hspace{-\tabcolsep}\enspace}ll
>{\raggedright\arraybackslash}X}
\lsptoprule
%\tbx{00-18}	&	lury	&	bone	&	< PMP *duRi thorn (bone in most of Wallacea)	\\	\midrule
%\endfirsthead
\tbx{13-18}	&	lury	&	bone	&	< PMP *duRi `thorn' (`bone' in most of Wallacea)	\\	\midrule
%\endhead
	&	lurkye	&	the bone	&	/lury-ke/  (\it{-ke} = \tsc{det})	\\
	&	lurikkwe	&	my (\tsc{1sg}) bone	&	/lury-V(\tsc{gen})-kw-ke/	\\
	&	lurikure	&	my (\tsc{1sg}) bones	&	/lury-V(\tsc{gen})-kw-Vre/ (-Vre = \tsc{pl})	\\
	&	lurimkwe	&	your (\tsc{2sg}) bone	&	/lury-V(\tsc{gen})-mw-ke/	\\
	&	lurimure	&	your (\tsc{2sg}) bones	&	/lury-V(\tsc{gen})-mw-Vre/	\\
	&	lurike	&	its (\tsc{3sg}) bone	&	/lury-V(\tsc{gen})-{\0}-ke/	\\
	&	lurinare	&	its (\tsc{3sg}) bones	&	/lury-V(\tsc{gen})-n-Vre/	\\
	&	ity luritare	&	our (\tsc{1pl.incl}) bones	&	/ity lury-V(\tsc{gen})-t-Vre/	\\
	&	ara lurimire	&	our (\tsc{1pl.excl}) bones	&	/ara lury-V(\tsc{gen})-my-Vre/	\\
\lspbottomrule
\end{tabularx}\mbox{}\\}

\citet[252]{Mills1991} also appeals to binding to explain the
morphophonemic consonant alternations in the inflectional paradigms
of languages in other parts of central Maluku (see \trf{07-16}
for Larike and Tables \ref{tab:07-17} and \ref{tab:12-01} for Banda)
for examples of such morphophenemic alternations).
\citet[252f]{Mills1991} notes that binding has also been observed
in Biak, Patani, Misool (SHWNG), and West Sumba, and adds “As 
Dyen (1978:248) [\citealp{Dyen1978a}]
says, binding ‘constitute[s] a reasonably strong argument for subgrouping
together with the languages exhibiting it’,
and I make just such use of it in positing the ‘SE Maluku’ subgroup.” \citep[253]{Mills1991}.
Some of what \citeauthor{Mills1991} identifies as binding are single instances
of what may or may not be a grammatical pattern for that language,
so more data are needed to establish that these actually are similar phenomena.

\citet[253]{Mills1991} goes on to reject borrowing
and parallel development as explanations for binding.
And then proceeds to speculate as follows:

\begin{quote}
		“The most attractive hypothesis would be:
		
		(c) All languages of
		Sumba, along with Proto Lettic,
		Proto Tanimbar-Kei, Numfor-Patani (and perhaps Proto Timor
		and Savu) ultimately descend from the same ‘Post-CMP’ stage;
		binding has been lost in most of Sumba
		and SHWNG (and all of Timor/Savu, if included).

		Thus we can see that it is not too difficult to create a reasonably
		coherent CMP subgroup, with binding, that embraces all the known descendants.
		The result, however, is to negate \citeauthor{Blust1978}’s (\citeyear{Blust1978}) argument that
		SHWNG and Oceanic are coordinate under a node labelled ‘Eastern Malayo-Polynesian’.
		Perhaps the situation could be salvaged by reformulating the
		SHWNG group without Numfor/Biak
		and Patani, though I suspect not.
		Blust’s arguments for this subgroup seem very strong,
		however one may feel about its alignment with Oceanic.

		More likely we should have to abandon our ‘strong assumption’
		of a single subgroup with binding,
		and allow after all for parallel development within a relatively
		restricted geographic area in north-eastern Indonesia.

		But if we reject parallel development,
		insisting on a single subgroup,
		and yet wish to retain Blust’s EMP hypothesis,
		then we are forced into one of these uncomfortable positions:

		(a) PTK [Proto-Tanimbar-Kei]/PLet [Proto-Lettic]/WSumba [West
		Sumba] as well as SHWNG,
		are \it{all} coordinate with OC [Oceanic],
		and binding is an innovation within EMP; or
		
		(b) PTK/etc. are indeed CMP languages
		and SHWNG is indeed EMP.
		Therefore, binding in both is \it{retained} from a higher level (Blust’s
		East-Central MP at least)
		and has been lost in a very large number of daughter languages
		(most CMP and apparently all OC [Oceanic]).
		
		I do not know the answer.”
		(\citealp[253]{Mills1991}; emphasis in the original)
\end{quote}

\citet[276]{Blust1993} rejects \citeauthor{Mills1991}’ use of ``binding'' as a basis
for justifying “Southeast Maluku” as a subgroup saying,
“However, ‘binding’ is not confined to the languages in which
Mills reports it, but is found in very similar form in such South Halmahera-West
New Guinea languages as Numfor (cf.
the appendix to Dyen (1978) [\citealp{Dyen1978a}],
where the subgrouping value of this innovation is discussed).
It thus seems premature to accept the existence of Mills’s ‘southeast
Maluku’ subgroup.” 
While we agree with Blust regarding the value of binding for
subgrouping “Southeast Maluku”,
largely because it ignores things we consider to be more basic
in the historical phonology,
nevertheless \citeauthor{Blust1993} has
not addressed the broader claims of \citeauthor{Mills1991}’
musings and their potential higher-level implications for subgrouping
in the wider region.

There are actually several different types of morpheme-internal,
morpheme-external, and other types of synchronic metathesis found
in languages of the region,
mostly in concentrated areas,
and some of which are triggered by high vowels.
\citet[20--76]{Edwards2020} provides a detailed survey of these
with data, discussion, and a map of languages in the greater Timor-Maluku
region that have different kinds of metathesis.
There is additional evidence for a considerable amount of diachronic
metathesis in the lexicons of the region,
as illustrated throughout this present study.\footnote{
	There are about 190 instances of the word \it{metathesis},
	or its abbreviations \it{(m.)/(met.)} in this study.
	While some of these refer to to the same forms,
	it illustrates the prevalence of metathesis in this region.}

The distribution of the occurrence or non-occurrence of binding,
along with other more basic differences in the historical phonologies
(see the relevant chapters),
suggests that binding is better treated as an areal phenomenon
from diffusion that crosses some subgrouping boundaries,
and is found in only some nodes of individual subgroups,
rather than as a genealogical feature that
forms the basis of a subgrouping argument.

It is also difficult to see how binding explains the wide variety
of morphophonemic alternations in subject prefixing in \tsc{Ambon-Seram}
languages, Banda, and Geser (the latter a \tsc{Seram-Tanimbar-Bomberai} language).
In those languages there is no other independent evidence to
support how the hypothetical underlying \it{i}-C or \it{u}-C
yields the great variety of consonants resulting from binding
(as opposed to the fairly complex but straight-forward labialization
and palatalization found in Yamdena,
Selaru, Luang, Leti, etc.).
This is particularly a problem since \it{i}
and \it{u} also maintain their identities in some inflectional
contexts in those languages.
We see them as different phenomena.
The morphophonemic alternations in \tsc{Ambon-Seram},
Banda and \tsc{Seram-Tanimbar-Bomberai} languages are described in \srf{07-04-03}.

For binding to be given more weight,
it would need to be reconstructed in certain morphemes for a
language ancestral to one or more of our subgroups.\footnote{
		Our thanks to Malcolm Ross for drawing this to our attention.}
Since it involves various morphophonemic processes relating to high
vowels and semi-vowels (or glides),
and different kinds of morphemes,
it is difficult to see how that might work.

Since binding is not found in all our subgroups,
and is also found outside the region (in some SHWNG languages,
but apparently not in Oceanic languages),
it is difficult to accept that it has value for subgrouping within
our target region, for either CMP or for CEMP,
to a degree that would override more basic patterns in the historical phonology.

Before we move on, it needs to be pointed out that \srf{04-06}
(see particularly example \qf{04-25} on \prf{ex:04-25} and surrounding discussion)
we \it{do} use some of the processes associated
with \citeauthor{Mills1991}' (\citeyear{Mills1991})
``binding'' for defining the \tsc{Southwest Maluku} node of \tsc{Timor-Babar}.
The evidence for \tsc{Southwest Maluku} comes from word-final VC{\#} {\ra} CV{\#} metathesis,
for which all languages of \tsc{Southwest Maluku} show some evidence.
While this metathesis does create words 
which are eligible to undergo ``binding''
(that is, V\tsc{[+high]}-final words), it is a separate process.
